\documentclass{ou}
\begin{document}
\thispagestyle{empty}
\begin{center}
\shortstack{\FGAVO}\vspace{4\baselineskip}
\textbf{РЕФЕРАТ}\\
\setstretch{2}
Информационные системы и~технологии\\
\setstretch{1}
Этапы жизненного цикла информационных систем\vspace{2\baselineskip}\end{center}
\begin{flushleft}
\doublespacing
\begin{tabularx}{\textwidth}{@{} p{0.365\textwidth} l l @{}}
Руководитель & ст. преподаватель & И. Р. Нигматуллин\\
Студенты     &                   &                  \\
\multirow{-2.27}{\hsize}{\singlespacing Понятие жизненного цикла информационной системы и~его значение}
             & ГФ25-01Б, 152537671 & А. Е. Корикова\\
             & ГФ25-01Б, 152536528 & Д. А. Васильева\\
             & ГФ25-01Б, 152540016 & И. Н. Морозова\\
\multirow{-2.27}{\hsize}{\singlespacing Анализ требований и~проектирование структуры будущей системы}
             & ГФ25-01Б, 152541294 & Э. Р. Махмудов\\
             & ГФ25-01Б, 152539829 & В. Д. Бобкова\\
\multirow{-2.27}{\hsize}{\singlespacing Этапы разработки и~внедрения информационной системы}
             & ГФ25-01Б, 152537684 & О. И. Бочанова\\
             & ГФ25-01Б, 152539986 & М. К. Гинтер\\
             & ГФ25-01Б, 152541471 & А. О. Безруких\\
             & ГФ25-01Б, 152536796 & А. В. Корюкова\\
\end{tabularx}
\end{flushleft}
\vspace*{\fill}
\begin{center}
\city
\end{center}
\newpage
\thispagestyle{empty}
\setstretch{1}
Продолжение титульного~листа реферата по~теме «Этапы жизненного цикла информационных систем»
\begin{flushleft}
\doublespacing
\begin{tabularx}{\textwidth}{@{} p{0.365\textwidth} l l @{}}
Студенты     &                     &         \\
\multirow{-2.27}{\hsize}{\singlespacing Тестирование, отладка и~интеграция функциональных компонентов}
             & ГФ25-01Б, 152540046 & А. Д. Балцевич\\
             & ГФ25-01Б, 152537659 & В. Е. Зайцева\\
             & ГФ25-01Б, 152536592 & А. Ю. Бирюкова\\
\multirow{-2.27}{\hsize}{\singlespacing Эксплуатация, сопровождение и~модернизация системы}
             & ГФ25-01Б, 152540753 & М. И. Бобылев\\
             & ГФ25-01Б, 152540765 & А. В. Гайворонская\\
             & ГФ25-01Б, 152539850 & Е. А. Аксаментова\\
\multirow{-2.27}{\hsize}{\singlespacing Завершение жизненного цикла и оценка эффективности информационной системы}
             & ГФ25-01Б, 152536115 & С. В. Идрисова\\
             & ГФ25-01Б, 152536056 & В. А. Карасова\\
             & ГФ25-01Б, 152539526 & К. Н. Аркадьева\\
             & ГФ25-01Б, 152540131 & Е. Кичко\\
\end{tabularx}
\end{flushleft}
\newpage
\tableofcontents
\newpage
\onehalfspacing
\section{Понятие жизненного цикла информационной системы и~его значение}
Информационные технологии и информационные системы оказали влияние на деятельность организаций по всему миру. Появление персональных компьютеров, создание локальных сетей, технологии «клиент-сервер» и сети Интернет позволило организациям нарастить возможности по контролю и~качеству управленческих решений в организационных структурах, вывело планирование финансовой и~экономической деятельности на более высокий качественный уровень. В связи с широкой компьютеризацией общества на~современном этапе предъявляются высокие требования к аппаратной части компьютеров, программному обеспечению и информационным системам.

Создание и функционирование информационной системы -- сложный и многоэтапный процесс, который требует четкой структуризации работ и определенной методологии внедрения. В связи с этим предлагается использование понятие жизненного цикла информационной системы, который представляет собой непрерывный процесс ее построения и развития технического задания вплоть до изъятия из эксплуатации. Каждый этап жизненного цикла включает в себя определенный состав, последовательность осуществляемых работ и их непосредственные результаты. Отдельный этап характеризуется различными методами и средствами, используемыми для~выполнения работ, а также различными ролями и ответственностью участников. Результатом такого детализированного описания этапов жизненного цикла служит четко спланированный и организованный процесс коллективной разработки информационной системы\cite{cyberleninka_lifecycle}.
\subsection{Жизненный цикл информационной системы и его структура}
Технологии проектирования, применяемые в настоящее время, предполагают поэтапную разработку системы. Этапы по общности целей могут объединяться в стадии.

Совокупность стадий и этапов, которые проходит ИС в своем развитии от момента принятия решения о создании системы до момента прекращения функционирования системы, называется жизненным циклом ИС.

Методология проектирования информационных систем описывает процесс создания и сопровождения систем в виде жизненного цикла (ЖЦ) ИС, представляя его как некоторую последовательность стадий (этапов) и~выполняемых на них процессов.

Для каждого этапа определяются:
\begin{itemize}[label=-, leftmargin=0em, itemindent=3.35em, nosep]
\item Состав и последовательность выполняемых работ;
\item Получаемые результаты;
\item Методы и средства, необходимые для выполнения работ;
\item Роли и ответственность участников и т.д.
\end{itemize}

Такое формальное описание ЖЦ ИС позволяет спланировать процесс коллективной разработки и обеспечить управление этим процессом.

Стадия -- часть процесса создания ИС, ограниченная определенными временными рамками и заканчивающаяся выпуском конкретного продукта (моделей, программных компонентов, документации), определяемого заданными для данной стадии требованиями.

Стадии жизненного цикла:
\begin{enumerate}[leftmargin=0em, itemindent=3.75em, nosep, label=\arabic*)]
\item Планирование и анализ требований -- системный анализ. Исследование и анализ существующей информационной системы, определение требований к создаваемой ИС, оформление технико-экономического обоснования и~технического задания на разработку ИС;
\item Проектирование. Разработка в соответствии со сформулированными требованиями состава автоматизируемых функций (функциональная архитектура) и состава обеспечивающих подсистем (системная архитектура), оформление технического проекта ИС;
\item Реализация (рабочее проектирование, физическое проектирование, программирование). Разработка и настройка программ, наполнение баз данных, создание рабочих инструкций для персонала, оформление рабочего проекта;
\item Внедрение (тестирование, опытная эксплуатация). Комплексная отладка подсистем ИС, обучение персонала, поэтапное внедрение ИС в~эксплуатацию, оформление акта о приемо-сдаточных испытаниях ИС;
\item Эксплуатация ИС (сопровождение, модернизация). Сбор рекламаций и статистики о функционировании ИС, исправление ошибок и недоработок, оформление требований к модернизации ИС и её выполнение \cite{isis_lifecycle}.
\end{enumerate}

Второй и третий этапы объединяют в одну стадию, называемую техно-рабочим проектированием или системным синтезом.

Завершается жизненный цикл информационной системы выводом ее из~эксплуатации.

Для каждой стадии определяют:
\begin{itemize}[label=-, leftmargin=0em, itemindent=3.35em, nosep]
\item Состав и последовательность выполняемых работ;
\item Получаемые результаты;
\item Методы и средства, необходимые для выполнения работ;
\item Роли и ответственность участников и т. д.
\end{itemize}

Такое формальное описание жизненного цикла информационной системы позволяет спланировать и организовать процесс коллективной разработки и обеспечить управление этим процессом. Рассмотрим каждую из~стадий более подробно.

Цели системного анализа:
\begin{enumerate}[leftmargin=0em, itemindent=3.75em, nosep, label=\arabic*)]
\item Сформулировать потребность в новой ИС (идентифицировать все недостатки существующей ИС);
\item Выбрать направление и определить экономическую целесообразность проектирования ИС\cite{pis_lecture}.
\end{enumerate}

Системный анализ ИС начинается с анализа функционирования объекта рассматриваемого экономического объекта (системы) в соответствии с~требованиями, которые предъявляются к нему. В результате этого этапа выявляется основные недостатки существующей ИС, на основе которых формулируется потребность в совершенствовании системы управления этим объектом и ставится задача определения экономически обоснованной необходимости автоматизации определенной функции управления, т.е.~создается технико-экономическое обоснование проекта.

После определения этой потребности появляется проблемах выбора направлений совершенствования объекта на основе выбора программных и~технических средств. Результаты оформляются в виде технического задания на проект, в котором определяется технические условия и требования к ИС, а~также ограничения на ресурсы проектирования.

Требования к ИС определяется в терминах функций, реализуемых системой, и предоставляемой ею информацией.

Системный синтез предполагает:
\begin{enumerate}[leftmargin=0em, itemindent=3.75em, nosep, label=\arabic*)]
\item Разработку функциональной архитектуры (ФА) ИС, которая отражает структуру выполняемых функций;
\item Разработку системной архитектуры (СА) выбранного варианта ИС, т.е.~состав обеспечивающих подсистем;
\item Выполнение реализации проекта \cite{cyberleninka_lifecycle}.
\end{enumerate}
 
Построение СА на основе ФА предполагает выделение элементов и модулей информационного, технического, программного обеспечения и др. обеспечивающих подсистем, определение связей по информации и~управлению между выделенными элементами и разработку технологии обработки информации.

Этап конструирования (физического проектирования) включает разработку инструкций пользователям и программ, создание ИО, включая наполнение баз данных.

Внедрение разработанного проекта заключается в опытном внедрении и промышленном внедрении. Опытное внедрение заключается в проверке работоспособности элементов и модулей проекта, устранении ошибок на~уровне элементов и связей между ними.

Сдача в промышленную эксплуатацию заключается в проверке проекта на~уровне функций и контроля соответствия его требованиям, сформулированным на стадии системного анализа.

Эксплуатация и сопровождение проекта включает в себя:
\begin{enumerate}[leftmargin=0em, itemindent=3.75em, nosep, label=\arabic*)]
\item Поддержка пользователей -- помощь в работе, консультации, устранение сбоев;
\item Контроль работы системы -- обеспечение стабильности, мониторинг производительности;
\item Внесение исправлений -- устранение ошибок, не выявленных при~тестировании;
\item Адаптация и развитие -- доработка системы под новые бизнес-задачи и изменения в законодательстве.
\end{enumerate}

Цель этапа -- обеспечить эффективное использование системы и ее постоянное соответствие потребностям бизнеса, которые были определены на~этапе системного анализа\cite{isis_lifecycle}.
\subsection{Модели жизненного цикла ИС}
В настоящее время известны и используются следующие модели ЖЦ:
\begin{itemize}[label=-, leftmargin=0em, itemindent=3.35em, nosep]
\item Каскадная модель (до 70-х годов);
\item Итерационная модель (70-80-ые годы);
\item Спиральная модель. (80 -- 90-ые годы)\cite{pis_lecture}.
\end{itemize}

Каскадная модель предусматривает последовательное выполнение всех этапов проекта в строго фиксированном порядке. Переход на следующий этап означает полное завершение работ на предыдущем этапе.

Для каскадной модели характерна автоматизация отдельных несвязанных задач, не требующая выполнения информационной интеграции программного, технического и организационного сопряжения.

Применение каскадной модели жизненного цикла к большим и сложным проектам вследствие большой длительности процесса проектирования и изменчивости требований за это время приводит к их практической нереализуемости.
 
Разработка ИС ведется итерациями с циклами обратной связи между этапами. Межэтапные корректировки позволяют учитывать реально существующее взаимовлияние результатов обработки на различных этапах; время жизни каждого из этапов растягивается на весь период разработки.

Создание комплексных ИС предполагает проведение увязки проектных решений, получаемых при реализации задач. Поход к проектированию «сверху-вниз» обуславливает необходимость таких итерационных возвратов, когда проектные решения по отдельным задачам комплектуются в общие системные решения и при этом возникает потребность в пересмотре ранее сформулированных требований. Вследствие большого числа итераций возникает рассогласования в выполненных проектных решениях.

На каждом витке спирали выполняется создание очередной версии продукта, уточняются требования проекта, определяется его качество и~планируются работы следующего витка.
 
На основе спиральной модели ЖЦ лежит применение прототипной технологии или RAD-технологии. ИС разрабатывается путем расширения программных прототипов повторяя путь от детализации требований к~детализации программного кода.

При прототипной технологии сокращается число итераций и меньше возникает ошибок и несоответствий, которые необходимо исправлять на~последующих итерациях, а само проектирование ИС осуществляется более быстрыми темпами, упрощается создание проектной документации \cite{cyberleninka_lifecycle}.
\subsection{Значение жизненного цикла ИС}
Понятие жизненного цикла информационной системы обеспечивает системный подход к созданию, внедрению и сопровождению сложных программно-аппаратных комплексов. Его значение проявляется в том, что он предоставляет комплексную методологическую основу для управления всем процессом разработки -- от первоначальной идеи до вывода системы из~эксплуатации\cite{isis_lifecycle}.

С методологической точки зрения жизненный цикл служит основой для~контроля проекта. Разбивая сложный процесс на~последовательные фазы, он позволяет реалистично оценивать сроки, бюджет и необходимые ресурсы. 

Особую значимость имеет то, что жизненный цикл обеспечивает соответствие создаваемой системы реальным потребностям бизнеса. На этапе системного анализа тщательно изучаются бизнес-процессы и формулируются требования, которые становятся ориентиром на всех последующих стадиях. 

Экономический аспект проявляется в оптимизации затрат на разработку и~сопровождение. Исследования показывают, что ошибка, обнаруженная на~этапе эксплуатации, может быть в 100 раз дороже в исправлении, чем если бы она была выявлена на стадии проектирования. Жизненный цикл смещает акцент на раннее выявление проблем, что существенно снижает совокупную стоимость владения системой\cite{pis_lecture}.
\newpage
\section{Анализ требований и~проектирование структуры будущей системы}
Процесс создания современной информационной системы представляет собой сложный комплекс взаимосвязанных этапов, среди которых ключевое значение имеют анализ требований и проектирование структуры. Качество выполнения этих начальных стадий жизненного цикла разработки влияет на~функциональную полноту, надежность, масштабируемость и способность системы удовлетворять потребности заказчика.
\subsection{Анализ требований к информационной системе}
Анализ требований является фундаментальным этапом, в ходе которого формируется понимание функционального назначения системы. Основная цель этого этапа заключается в полном, непротиворечивом и документированном определении функциональных и нефункциональных потребностей заказчика и~будущих пользователей системы.

Процесс анализа требований включает выполнение нескольких подэтапов. На этапе выявления и сбора требований применяются методы получения информации, такие как интервьюирование заинтересованных лиц, анализ документации и бизнес-процессов организации, проведение опросов, наблюдение за работой пользователей в операционной среде, а также работа с~фокус-группами для проработки конкретных сценариев.

Анализ и систематизация требований предполагает проверку собранной информации на предмет противоречивости, полноты и~выполнимости. На~этом этапе требования классифицируются на функциональные, определяющие конкретные функции и операции системы, и нефункциональные, описывающие общие характеристики системы и~ограничения ее работы.

Завершающим подэтапом процесса является спецификация требований, результатом которой становится формирование официального документа -- Технического Задания или спецификации требований к программному обеспечению. Этот документ служит основным источником информации для~всех участников проекта и юридической основой для приемки системы.
\subsection{Проектирование структуры информационной системы}
После четкого определения и согласования требований наступает этап проектирования, который отвечает на вопрос о способах реализации системы. Проектирование структуры направлено на создание надежной, эффективной и~поддерживаемой архитектуры системы.

Процесс проектирования структуры охватывает несколько аспектов. Архитектурное проектирование предполагает принятие стратегических решений об общей структуре информационной системы, включая выбор архитектурного стиля, определение основных компонентов системы и их взаимодействия, а также распределение функций по логическим уровням.

Проектирование данных фокусируется на структурах хранения и~управления информацией, включая разработку логической и физической модели данных, проектирование схемы базы данных, нормализацию и~планирование индексов для оптимизации выполнения запросов.

Детальное проектирование компонентов предполагает их проектирование с высокой степенью детализации, включая определение интерфейсов компонентов, описание алгоритмов и бизнес-логики, а также спецификацию форматов данных для обмена между компонентами.

Важным аспектом является проектирование пользовательского интерфейса и взаимодействия, в рамках которого разрабатываются макеты экранов, навигационные схемы и сценарии взаимодействия пользователя с~системой.

Результат этапа проектирования -- пакет проектной документации, содержащий архитектурные диаграммы, схемы базы данных, спецификации API и прототипы интерфейсов.

Анализ требований и проектирование структуры представляют собой два взаимосвязанных и критически важных этапа создания информационной системы. Анализ требований формирует ясное видение целей и задач системы, фиксируя их в виде формализованных требований. Проектирование структуры трансформирует эти требования в технически обоснованный план реализации, определяющий архитектуру, компоненты, данные и интерфейсы будущей информационной системы.
\newpage
\section{Этапы разработки и~внедрения информационной системы}
Разработка и внедрение информационной системы (ИС) -- это сложный и многогранный процесс, направленный на автоматизацию и оптимизацию бизнес-процессов организации, а также на повышение эффективности управления и взаимодействия внутри предприятия и с внешней средой. В~современном мире, где конкуренция нередко определяется скоростью и~обработки информации, правильное и своевременное внедрение ИС становится ключевым фактором успеха.

Особенности этого процесса определяются методологией проектирования системы -- строительно-структурной, объектно-ориентированной. Важными аспектами являются масштабы организации, отраслевые требования и~специфические особенности, связанные с корпоративной культурой. Внедрение системы предполагает не только техническую реализацию, но~и~значительную социальную адаптацию, обучение персонала и изменение внутренних процессов. В результате правильное планирование, подготовка и проведение каждого этапа позволяют минимизировать риски, такие как сопротивление изменениям персонала или чрезмерные затраты.
\subsection{Основные этапы разработки и внедрения}
Анализировав разнообразие подходов, можно выделить основные этапы, универсальные для большинства проектов, несмотря на их специфику. Эти этапы могут реализовываться по модели, когда последовательно идут этапы:
\begin{enumerate}[leftmargin=0em, itemindent=3.75em, nosep, label=\arabic*)]
\item Проектирование;
\item Программирование;
\item Тестирование;
\item Внедрение.
\end{enumerate}

Либо по итеративным моделям, предусматривающим повторение и~уточнение каждого этапа в процессе реализации. Важную роль в управлении этими процессами играют современные средства моделирования, такие как IDEF0 (SADT), BPMN или DFD, которые позволяют визуализировать бизнес-процессы и структурировать требования.
\subsection{Анализ и планирование}
Первым и ключевым этапом является анализ предметной области организации и формирование требований к будущей системе. Этот шаг включает проведение технико-экономического анализа, который позволяет понять возможности предприятия, определить его сильные и слабые стороны, а также выявить потребности в автоматизации. Важна оценка структурных особенностей организации, корпоративной культуры, а также первичный анализ бизнес-процессов, которых касается внедрение ИС. Создается модель текущих процессов («как есть») и формируются целевые модели («как должно быть»), что служит основой для последующих этапов.

В авиационной отрасли при анализе системы особое внимание уделяется интеграции с существующими информационными ресурсами, такими как системы ERP или системы управления документооборотом (например, Directum). В многоуровневых и крупномасштабных организациях критически важным является тщательное изучение нормативных требований и особенностей взаимодействия различных подразделений, что обеспечивает более точное определение масштабов и этапов проекта.
\subsection{Проектирование и моделирование}
На этапе проектирования разрабатывается архитектура системы, ее логическая и физическая модели. Используются методы функционального моделирования, такие как SADT или BPMN, для структурирования бизнес-процессов и определения функционального состава системы. Важной задачей является формирование технического задания, в котором учитываются нефункциональные требования: производительность, безопасность.

Проектирование включает создание схем взаимодействия межу программными модулями, описание интерфейсов, а также определение технологического стека. Для автомобильного бизнеса важным аспектом становится разработка интерфейсов порталов с возможностью разграничения доступа по различным ролям и группам пользователей, что повышает безопасность и удобство эксплуатации системы. В транспортных организациях разрабатываются структуры автоматизированных рабочих мест (АРМ), обеспечивающих единое управление данными.
\subsection{Разработка системы}
После завершения этапа проектирования начинается непосредственная разработка программного обеспечения. В рамках этого этапа создаются программные модули, базы данных и осуществляется их интеграция с внешними системами и сервисами. Важную роль играет настройка инфраструктуры, инсталляция программных продуктов и организация безопасной среды работы.

В разработке используют декомпозицию процессов, основанную на~моделях SADT или IDEF3, что позволяет обеспечить четкое соответствие программных решений требованиям, вытекающим из модели «как должно быть». Например, в системах ERP (SAP, Oracle) осуществляется моделирование бизнес-процессов для настройки модулей и их адаптации под специфику предприятия.
\subsection{Тестирование и отладка}
На этом этапе проверяют готовую систему на предмет соответствия техническому заданию и требованиям, выявляют и устраняют возможные ошибки. Тестирование включает функциональную проверку, нагрузочные тесты и оценку устойчивости системы к возможным сбоям. В большинстве проектов важной составляющей является обучение пользователей и подготовка технического персонала к эксплуатации системы.

Примером является портальная система на базе 1С-Битрикс, где этап приемочного тестирования может занимать до 270 дней. В автомобильной промышленности после успешного тестирования осуществляется подготовка к массовому внедрению и началу эксплуатации, что включает обучение сотрудников и настройку всех элементов системы.
\subsection{Внедрение системы}
Развертывание системы предполагает введение в~эксплуатацию, что~особенно актуально для сложных крупномасштабных проектов. В~зависимости от стратегии внедрения, применяются разные подходы: полномасштабный «большой взрыв», поэтапное внедрение или франчайзинг. Важной задачей на этом этапе является миграция данных, обучение сотрудников, изменение бизнес-процессов под условия новой системы. В~автомобильной отрасли это может означать подключение дилеров, сервисных центров и обеспечение обмена информацией в реальном времени.

Общий срок внедрения крупных систем нередко превышает 400 дней, что обусловлено необходимостью масштабных настроек, развертывания инфраструктуры и обучения персонала. В этой фазе важно обеспечить постоянную обратную связь, мониторинг работы системы и адаптацию процессов под возникающие новые требования.
\subsection{Сопровождение и эксплуатация}
После внедрения система требует постоянной поддержки. В рамках этого этапа ведется мониторинг ее работы, устраняются возникающие сбои и~реализуются обновления. Важной задачей является формирование регламентов эксплуатации, закрепление правил работы с системой и~мотивация сотрудников к использованию новых инструментов. В условиях строго регулируемой среды, например, в сфере финансирования или авиации, особое внимание уделяется вопросам безопасности данных и соответствию нормативным актам (например, ФЗ №152).

Постоянное совершенствование системы включает адаптацию к~изменениям бизнес-стратегии, внедрение новых технологий и расширение функциональности, что способствует поддержанию конкурентоспособности организации.

Этапы разработки и внедрения информационной системы связаны между собой и образуют цикличный, итеративный процесс, результат которого зависит от правильного выбора методологии, участия коллектива и внимания ко всем аспектам проектирования и реализации. В условиях современной динамичной среды успешное внедрение ИС требует тщательного предварительного анализа, четкого планирования и профессиональной реализации каждого этапа. Ключевые вызовы -- человеческий фактор, риски превышения бюджета и сложности интеграции -- требуют системного подхода и привлечения консультантов, а также предпроектных исследований для~уменьшения неопределенности и повышения вероятности успеха.
\newpage
\section{Тестирование, отладка и интеграция функциональных компонентов}
\subsection{История тестирования}
Тестирование программ зародилось практически одновременно с~программированием. Машинное время стоило дорого, поэтому предприятия с~повышенными требованиями к надежности программ стали активно разрабатывать методики тестирования.

Долгое время было принято считать, что целью тестирования является доказательство отсутствия ошибок в программе. Однако этот тезис не выдерживает критики, т. к. полный перебор всех возможных вариантов выполнения программы находится за пределами вычислительных возможностей даже для очень небольших программ. 
\subsection{Тестирование: общие сведения}
Определение для тестирования предложил один из его основоположников Гленфорд Майерс: «Тестирование -- это процесс выполнения программ с целью обнаружения ошибок». Здесь следует отметить, что тестовая деятельность, предусматривающая эксплуатацию программного продукта, называется динамическим тестированием (англ., Dynamic Testing). В то~же~время существует понятие статического тестирования (англ., Static Testing), определяющего тестовую деятельность, связанную с анализом программного обеспечения (ПО) и проводимую без его запуска. К статическому тестированию относятся обзоры кода, инспекции, аудит, критический анализ и т. д. Внимательное изучение этих двух подходов к тестированию показывает, что они дополняют друг друга и позволяют находить разные виды ошибок. Поэтому наиболее эффективные процессы разработки ПО используют комбинацию методов тестирования, реализующих эти подходы. 

Тестирование максимально эффективно в тех случаях, когда программа проверяется не только путем запуска, но и путем чтения, статических проверок и т. п. Поэтому классическое определение Майерса, приведенное выше, оказывается слишком узким и не охватывающим всех аспектов современного тестирования. Поэтому будем пользоваться следующим определением: Тестирование ПО (англ., Software Testing) -- это процесс анализа или эксплуатации программного обеспечения с целью выявления дефектов. Слово «процесс» используется для уточнения, что тестирование -- это плановая и упорядоченная деятельность. 

Цель тестирования -- выявление ошибок, недочетов и несоответствий спецификациям до того, как продукт будет выпущен на рынок.

Тестирование можно классифицировать по очень большому количеству признаков. Вот некоторые из них:
\begin{itemize}[label=-, leftmargin=0em, itemindent=3.35em, nosep]
\item По запуску кода на исполнение;
\item По доступу к коду и архитектуре приложения;
\item По степени автоматизации;
\item По уровню детализации приложения (по уровню тестирования);
\item По (убыванию) степени важности тестируемых функций (по уровню функционального тестирования);
\item По принципам работы с приложением;
\end{itemize}

По запуску кода на исполнение статическое тестирование без запуска и~динамическое тестирование с запуском.

По доступу к коду и архитектуре приложения метод белого ящика -- изучение внутреннего устройства программы (исходных текстов, модулей и~их структур), метод чёрного ящика -- программе подаются некоторые данные на~вход и проверяются результаты в надежде найти несоответствия, при этом как именно работает программа считается несущественным, и~этот подход до~сих пор является самым распространённым в повседневной практике, а~метод серого ящика -- имеется частичный доступ к исходному коду (например, сторонние компоненты или автоматизация функционального тестирования через пользовательский интерфейс).

По степени автоматизации ручное тестирование -- тест-кейсы выполняет человек, и автоматизированное тестирование -- тест-кейсы частично или~полностью выполняет специальное инструментальное средство.

По уровню детализации приложения (по уровню тестирования) модульное (компонентное) тестирование -- проверяются небольшие части приложения, интеграционное тестирование -- проверяется взаимодействие между несколькими частями приложения, и системное тестирование -- приложение проверяется как единое целое.

По (убыванию) степени важности тестируемых функций (по уровню функционального тестирования) дымовое тестирование -- проверка самой важной, самой ключевой функциональности, неработоспособность которой делает бессмысленной саму идею использования приложения, тестирование критического пути -- проверка функциональности, используемой типичными пользователями в типичной повседневной деятельности, и расширенное тестирование -- проверка функциональности, заявленной в требованиях.

По принципам работы с приложением позитивное тестирование -- все действия с приложением выполняются строго по инструкции без никаких недопустимых действий, некорректных данных и т. д., можно образно сказать, что приложение исследуется в «тепличных условиях», и негативное тестирование -- в работе с приложением выполняются (некорректные) операции и используются данные, потенциально приводящие к ошибкам.

Тестирование может включать в себя составление тест-плана. Тест план (Test Plan) -- это документ, описывающий весь объем работ по~тестированию, начиная с описания объекта, стратегии, расписания, критериев начала и~окончания тестирования, до необходимого в процессе работы 10 оборудования, специальных знаний, а также оценки рисков с вариантами их разрешения.

Инструменты тестирования:
\begin{enumerate}[leftmargin=0em, itemindent=3.75em, nosep, label=\arabic*)]
\item Selenium: Один из самых популярных инструментов для автоматизации тестирования веб-приложений. Selenium поддерживает множество языков программирования, включая Java, C\verb|#|, Python;
\item JUnit и TestNG: Фреймворки для модульного тестирования, написанные на Java. TestNG считается более гибким и мощным благодаря более широкому спектру тестовых конфигураций и опций;
\item QTP/UFT (Unified Functional Testing): Коммерческий инструмент от~HP для функционального тестирования. Позволяет выполнять тестирование как десктопных, так и веб-приложений;
\item Cucumber -- это инструмент, поддерживающий Behavior-Driven Development. Позволяет писать тестовые сценарии на естественном языке, что~облегчает коммуникацию между разработчиками, QA;
\item Postman: Популярный инструмент для тестирования API. Позволяет удобно отправлять запросы к API, проверять ответы и автоматизировать тесты;
\item LoadRunner и JMeter: инструменты для нагрузочного тестирования. Позволяют проверить способность приложения выдерживать большое количество одновременных пользовательских запросов.
\end{enumerate}
\subsection{Об экономической стороне тестирования}
Тестирование является затратной деятельностью, отнимающей время и деньги. Поэтому в большинстве случаев разработчики программного обеспечения заранее формулируют какой-либо критерий качества создаваемых программ (определяют так называемую планку качества), добиваются выполнения этого критерия и затем прекращают тестирование, выпуская продукт на рынок. Такая концепция получила название Good Enough Quality (достаточно хорошее ПО), в противовес более очевидной концепции Best Possible Quality (максимально качественное ПО).

К сожалению, принцип Good Enough Quality зачастую понимают неправильно, ближе к формулировке Quality -- If Time Permits (качество, если будет время). Конечно, выпуск плохо протестированного программного продукта из-за недостатка времени -- это плохая практика. Опыт показывает, что пользователи склонны со временем забывать даже значительные задержки с~выпуском продукта, но плохое качество выпущенного продукта запоминается на всю жизнь. На самом деле, Good Enough Quality -- это просто поиск разумного компромисса между затратами на тестирование, длительностью разработки продукта и его качеством.
\subsection{Психологические аспекты тестирования}
Тестирование и программирование отличаются коренным образом, и~для~тестирования требуется иной склад характера, чем для программирования. Следовательно, программист не должен тестировать свои собственные программы -- для этого необходим другой человек. Более того, программист не должен тестировать даже чужие про граммы! Дело в том, что программист потратит какое-то время, перенастраиваясь на новую задачу. Даже переключение с одной программистской задачи на другую требует обычно от 10 минут до~получаса. Поэтому переключение с одного образа мышления на другой отнимет существенно больше времени.
\subsection{Отладка}
Отладка -- это процесс обнаружения причин возникновения ошибок и их последующего исправления.

В методике отладки принято выделять две части:
\begin{itemize}[label=-, leftmargin=0em, itemindent=3.35em, nosep]
\item Нахождение причины возникновения ошибки (90 \verb|%| времени отладки, как правило, уходит на эту часть); 
\item Исправление причины ошибки.
\end{itemize}

Наиболее эффективным считается обобщенный научный подход нахождения причин ошибок. Он состоит из следующих шагов:
\begin{enumerate}[leftmargin=0em, itemindent=3.75em, nosep, label=\arabic*)]
\item Сбор данных через повторяющиеся эксперименты;
\item Создание гипотезы, отражающей максимум доступных данных;
\item Разработка эксперимента для проверки гипотезы; 
\item Подтверждение или опровержение гипотезы;
\item Итерация предыдущих шагов (при необходимости).
\end{enumerate}

Тогда, с учетом обобщенного научного подхода, методику отладки можно описать в виде последовательности следующих этапов:
\begin{itemize}[label=-, leftmargin=0em, itemindent=3.35em, nosep]
\item Стабилизация ошибки (1-й шаг обобщенного научного подхода);
\item Обнаружение точного места ошибки (2--4 шаги обобщенного научного подхода);
\item Исправление ошибки;
\item Проверка исправления;
\item Поиск похожих ошибок.
\end{itemize}
\subsection{Инструменты отладки}
Отладчики позволяют разработчикам исполнять программу пошагово, останавливая её выполнение на определенных точках (брейкпоинты), чтобы~изучить текущее состояние программы. Это включает в себя возможность просмотра значений переменных, состояние стека вызовов и поток выполнения программы. Популярные отладчики:
\begin{itemize}[label=-, leftmargin=0em, itemindent=3.35em, nosep]
\item GDB для программ на C/C++ в Unix/Linux-системах;
\item PDB для Python-программ;
\item Visual Studio Debugger для .NET и C++ проектов в среде Windows;
\item LLDB для проектов на C/C++ и Swift поддерживающих Clang.
\end{itemize}

Логирование -- это процесс записи событий, происходящих в программе, во внешние файлы или системы хранения логов. Логи могут содержать информацию о работе программы, предупреждениях и ошибках. Логирование полезно для отслеживания поведения программы в реальном времени и~для~постфактум-анализа проблем. Инструменты для логирования:
\begin{enumerate}[leftmargin=0em, itemindent=3.75em, nosep, label=\arabic*)]
\item Log4j для Java-приложений;
\item NLog и Serilog для .NET-проектов;
\item Winston и Morgan для Node.js-приложений.
\end{enumerate}

Диагностические инструменты помогают анализировать программу для~выявления причин проблем, не всегда непосредственно связанных с~ошибками в коде, например, утечки памяти. Примеры таких инструментов:
\begin{itemize}[label=-, leftmargin=0em, itemindent=3.35em, nosep]
\item Valgrind для обнаружения некорректного использования памяти в~программах на C/C++;
\item Profiler в Visual Studio для анализа производительности и использования ресурсов .NET приложений;
\item Chrome Developer Tools для веб-разработки, включающие инструменты для отладки JavaScript, анализа производительности веб-страниц и многое другое.
\end{itemize}

Типы ошибок:
\begin{enumerate}[leftmargin=0em, itemindent=3.75em, nosep, label=\arabic*)]
\item Синтаксические ошибки -- ошибки, вызванные нарушением правописания или пунктуации в записи выражений, операторов и т. п., иначе говоря, нарушении грамматических правил языка. Обнаруживаются компилятором;
\item Семантические ошибки -- нарушение порядка операторов, параметров функций и употреблении выражений. Семантические ошибки также обнаруживаются компиляторо;
\item Логические ошибки -- труднее всего обнаружить. Они не вызывают сбой программы и не блокируют еѐ запуск, они заставляют ее каким-то образом «плохо себя вести», отображая неправильный вывод. Компилятор может выявить только следствие логической ошибки.
\end{enumerate}

Методы отладки (на основе логики):
\begin{itemize}[label=-, leftmargin=0em, itemindent=3.35em, nosep]
\item Ручное тестирование;
\item Метод индукции;
\item Дедукция;
\item Метод обратного прослеживания.
\end{itemize}

Ручное тестирование -- самый простой и естественный способ. При~обнаружении ошибки крайне важно выполнить тестируемую программу вручную, используя тестовый набор, при работе с которым была обнаружена ошибка.
 
Метод очень эффективен, но не применим для больших программ, программ со сложными вычислениями и в тех случаях, когда ошибка связана с~неверным представлением программиста о выполнении некоторых операций.
 
Метод индукции основан на тщательном анализе симптомов ошибки, которые могут проявляться как неверные результаты вычислений или как сообщение об ошибке. В случае если компьютер просто «зависает», то~фрагмент проявления ошибки вычисляют, исходя из последних полученных результатов и действий пользователя. Полученную таким образом информацию организуют и тщательно изучают, просматривая соответствующий фрагмент программы. В результате этих действий выдвигают гипотезы об ошибках, каждую из которых проверяют. В случае если гипотеза верна, то детализируют информацию об ошибке, иначе -- выдвигают другую гипотезу.
 
По методу дедукции сначала формируют множество причин, которые могли бы вызвать данное проявление ошибки. Затем, анализируя причины, исключают, те, которые противоречат имеющимся данным. В случае если все причины исключены, то следует выполнить дополнительное тестирование исследуемого фрагмента. В противном случае наиболее вероятную гипотезу пытаются доказать. В случае если гипотеза объясняет полученные признаки ошибки, то ошибка найдена, иначе -- проверяют следующую причину.
 
Метод обратного прослеживания наиболее эффективен для небольших программ. В данном случае отладку начинают с точки вывода неправильного результата. Для этой точки строится гипотеза о значениях базовых переменных, которые могли бы привести к получению имеющегося результата. Далее, исходя из этой гипотезы, делают предположения о значениях переменных в~предыдущей точке.

Методы отладки (статические):
\begin{itemize}[label=-, leftmargin=0em, itemindent=3.35em, nosep]
\item Ручной анализ кода (code review);
\item Статический анализ (static analysis);
\item Анализ структуры и архитектуры.
\end{itemize}

При ручном анализе кода программисты читают код, проверяют алгоритмы, корректность условий и структуру функций.

В статическом анализе специальные инструменты автоматически находят проблемы утечки памяти, неинициализированные переменные или потенциальные null-pointer ошибки. Примеры инструментов: PVS-Studio, SonarQube, ESLint (для JS).

Анализ структуры и архитектуры -- поиск логических конфликтов (неправильные связи между модулями, избыточная сложность, нарушение принципов SOLID)

Методы отладки (динамические):
\begin{enumerate}[leftmargin=0em, itemindent=3.75em, nosep, label=\arabic*)]
\item Отладка с пошаговым выполнением - используется в IDE: Visual Studio, PyCharm, VS Code, Xcode. Функции: breakpoints (точки остановки), step into / step over / step out, просмотр переменных, изменение значений переменных во~время выполнения;
\item Журналирование (logging) -- вывод событий. Журналирование используется в~веб-сервисах, мобильных приложениях, микросервисах;
\item Трассировка (tracing) -- следование за выполнением кода -- вывод шагов и состояний;
\item Использование динамических анализаторов -- инструменты, отслеживающие ошибки в процессе работы.
\end{enumerate}
\subsection{Интеграция}
Интеграция функциональных компонентов -- это процесс объединения отдельных частей системы (модулей, сервисов, программных блоков) в единую работающую систему, где все компоненты взаимодействуют корректно, обмениваются данными и выполняют общие функции.

Задачи интеграции:
\begin{itemize}[label=-, leftmargin=0em, itemindent=3.35em, nosep]
\item Компоненты создаются разными командами и нужно правильно их соединить;
\item Система должна быть согласованной, надёжной и функционировать как единое целое.
\end{itemize}

Способы интеграции компонентов:
\begin{enumerate}[leftmargin=0em, itemindent=3.75em, nosep, label=\arabic*)]
\item Горизонтальная интеграция -- при ней компоненты одного уровня взаимодействуют напрямую; 
\item Вертикальная интеграция -- интеграция создаёт «столбец» приложения: интерфейс -- логика -- данные -- инфраструктура;
\item Интеграция через шину (ESB -- Enterprise Service Bus) -- все компоненты подключаются через единую коммуникационную шину. Используется в больших корпоративных системах;
\item Интеграция через API -- стандартный современный подход;
\item Интеграция микросервисов,
\end{enumerate}

Микросервисы связаны через:
\begin{itemize}[label=-, leftmargin=0em, itemindent=3.35em, nosep]
\item HTTP-запросы;
\item Очереди сообщений (RabbitMQ, Kafka);
\item Брокеры событий.
\end{itemize}

Проблемы, которые решает интеграция:
\begin{enumerate}[leftmargin=0em, itemindent=3.75em, nosep, label=\arabic*)]
\item Несоответствие форматов данных;
\item Несовпадение интерфейсов и протоколов;
\item Ошибки синхронизации.
\end{enumerate}
\newpage
\section{Эксплуатация, сопровождение и~модернизация системы}
\subsection{Понятие эксплуатации информационных систем}
Эксплуатация информационной системы -- это совокупность действий, направленных на поддержание работоспособности ИС на протяжении всего жизненного цикла. В эксплуатацию включаются технические, организационные и технологические мероприятия, обеспечивающие корректное функционирование программного и аппаратного обеспечения.

Основные задачи эксплуатации:
\begin{enumerate}[leftmargin=0em, itemindent=3.75em, nosep, label=\arabic*)]
\item Обеспечение бесперебойной работы ИС. Это включает контроль состояния оборудования, операционных систем, сетевой инфраструктуры и~прикладного ПО;
\item Управление доступом и безопасностью. В эксплуатационной фазе используются политики парольной защиты, многофакторная аутентификация, контроль ролей пользователей;
\item Резервное копирование и восстановление данных. Важная задача эксплуатации -- организация бэкапов, проверка их целостности и готовности к восстановлению. Мониторинг и регистрация событий. Отслеживаются лог-файлы, нагрузка, сетевые взаимодействия.
\end{enumerate}

Пример эксплуатации: в производственной компании внедрена система учёта складских запасов. В рамках эксплуатации специалисты ИТ-отдела ежедневно контролируют корректность записи данных сканерами штрих-кодов, проводят проверку актуальности базы данных и обеспечивают стабильное подключение мобильных терминалов к Wi-Fi сети. Раз в неделю выполняется автоматическое резервное копирование на внешний сервер, расположенный в~другом офисе.

Cопровождение информационной системы -- деятельность по обновлению и совершенствованию системы в процессе эксплуатации без значительных изменений её архитектуры.

Виды сопровождения:
\begin{enumerate}[leftmargin=0em, itemindent=3.75em, nosep, label=\arabic*)]
\item Корректирующее сопровождение. Исправление ошибок, возникших в процессе эксплуатации: логических ошибок в коде, сбоев в интерфейсе, некорректного отображения данных;
\item Адаптивное сопровождение. Настройка системы под изменения внешних условий: обновления операционных систем, изменение тарифов и~ставок;
\item Совершенствующее сопровождение. Добавление незначительных функций, оптимизация производительности, улучшение интерфейса;
\item Предупреждающее сопровождение. Превентивные меры, нужные для~предотвращения будущих сбоев: обновление безопасности, оптимизация базы данных, установление новых защитных механизмов.
\end{enumerate}

Пример сопровождения: государственная организация использует информационную систему для электронного документооборота. После изменения закона о персональных данных ИТ-служба выполнила адаптивное сопровождение: обновила модули защиты информации, внедрила новые шаблоны согласования документов и изменила структуру хранилища, чтобы~обеспечить соответствие новым юридическим нормам.
\subsection{Модернизация информационных систем}
Модернизация -- это процесс обновления информационной системы, направленный на значительное улучшение её характеристик, расширение возможностей или повышение эффективности. В отличие от сопровождения, модернизация часто затрагивает архитектуру системы, включает замену оборудования или переход на новые программные решения.

Цели модернизации:
\begin{itemize}[label=-, leftmargin=0em, itemindent=3.35em, nosep]
\item Повышение производительности. Модернизация позволяет сократить время обработки данных, увеличить пропускную способность сервера или~системы хранения;
\item Поддержка новых технологий. Например, внедрение облачной инфраструктуры или использование искусственного интеллекта для анализа данных;
\item Оптимизация затрат. Обновление инфраструктуры может снизить стоимость её обслуживания в долгосрочной перспективе;
\item Повышение надежности и безопасности. Замена устаревших компонентов снижает риск отказов и уязвимостей.
\end{itemize}

Методы модернизации:
\begin{enumerate}[leftmargin=0em, itemindent=3.75em, nosep, label=\arabic*)]
\item Техническая модернизация. Замена серверов, систем хранения данных, сетевых маршрутизаторов;
\item Программная модернизация. Переход на новые версии программ, внедрение модульных архитектур, перенос части системы в облако;
\item Архитектурная модернизация. Полная реконструкция системы, переход от монолита к микросервисам;
\item Функциональная модернизация. Добавление новых крупных модулей или интеграция с другими системами.
\end{enumerate}

Пример модернизации: розничная сеть, использующая устаревшую систему кассового обслуживания, столкнулась с проблемой низкой скорости обработки транзакций и высокой нагрузкой в часы пик. Компания провела модернизацию: перевела базу данных на распределённую архитектуру, обновила POS-терминалы и внедрила облачный модуль аналитики. В~результате время обслуживания клиентов сократилось на 30\verb|%|, а стабильность системы значительно выросла.
\subsection{Жизненный цикл ИС и роль эксплуатации, сопровождения и модернизации}
Эксплуатация, сопровождение и модернизация -- это ключевые этапы жизненного цикла информационной системы. Система, созданная и внедрённая без продуманной стратегии сопровождения или модернизации, быстро устаревает и становится неэффективной.

Взаимосвязь процессов:
\begin{enumerate}[leftmargin=0em, itemindent=3.75em, nosep, label=\arabic*)]
\item Эксплуатация обеспечивает стабильную работу ИС на ежедневной основе;
\item Сопровождение помогает поддерживать систему актуальной и~функциональной;
\item Модернизация обеспечивает обновление ИС в стратегической перспективе, адаптируя её к новым условиям и требованиям.
\end{enumerate}

Пример комплексного подхода: в банке действует централизованная ИС управления кредитами.

На этапе эксплуатации поддерживается работа серверов, осуществляется резервное копирование, мониторинг активности пользователей. В процессе сопровождения постоянно обновляются справочники процентных ставок, добавляются новые формы отчётности. Проводится модернизация -- например, внедрение новых модулей скоринга на основе машинного обучения.
\subsection{Проблемы и риски при эксплуатации и модернизации}
Любая работа с информационными системами связана с определёнными рисками:
\begin{itemize}[label=-, leftmargin=0em, itemindent=3.35em, nosep]
\item Человеческий фактор. Ошибки администраторов, невнимательность пользователей;
\item Устаревшее оборудование. Приводит к сбоям, медленной работе и~высоким затратам;
\item Несовместимость ПО. Возникает при модернизации или обновлении модулей;
\item Киберугрозы. Вредоносное ПО, фишинг-атаки, попытки взлома;
\item Недостаточная документация. Это усложняет сопровождение и~модернизацию. Для минимизации рисков необходимо использовать стандарты ITIL, ISO/IEC 20000, а также вести строгую документацию всех изменений.
\end{itemize}
\newpage
\section{Завершение жизненного цикла и оценка эффективности информационной системы}
Жизненный цикл информационной системы (ИС) -- это период времени от момента возникновения потребности в системе до ее полного вывода из~эксплуатации. Завершение жизненного цикла ИС -- это неизбежная стадия, которая наступает вследствие морального или физического устаревания системы, изменения бизнес-требований или появления более эффективных альтернатив. Правильное и планомерное завершение жизненного цикла ИС так~же важно, как и её разработка, поскольку позволяет минимизировать риски и издержки, связанные с устаревшей системой. Важной составляющей процесса завершения является оценка эффективности ИС, позволяющая учесть опыт эксплуатации и принять обоснованное решение о дальнейшей судьбе системы.
\subsection{Причины завершения жизненного цикла ИС}
Завершение жизненного цикла ИС может быть обусловлено следующими основными причинами:
\begin{enumerate}[leftmargin=0em, itemindent=3.75em, nosep, label=\arabic*)]
\item Технологическое устаревание: использование устаревших технологий, которые сложно поддерживать, модернизировать и интегрировать с новыми системами. Отсутствие квалифицированных специалистов для поддержки и~развития устаревших технологий;
\item Функциональное устаревание: несоответствие функциональности системы изменяющимся бизнес-требованиям и новым стратегиям организации. Необходимость в новых функциях и возможностях, которые сложно или~невозможно реализовать в текущей системе;
\item Экономическое устаревание: высокие затраты на поддержание и эксплуатацию системы, превышающие экономическую выгоду от ее использования. Наличие более эффективных и экономически выгодных альтернатив;
\item Нормативное устаревание: несоответствие новым требованиям, стандартам и законодательству. Необходимость в приведении системы в~соответствие с новыми требованиями, что может потребовать существенных изменений или полной замены системы;
\item Физическое устаревание: износ оборудования, выход из строя компонентов системы, что приводит к снижению надежности и увеличению затрат на ремонт.
\end{enumerate}
\subsection{Процесс завершения жизненного цикла ИС}
Процесс завершения жизненного цикла ИС включает в себя ряд этапов, направленных на планомерный вывод системы из эксплуатации с~минимальными рисками:
\begin{itemize}[label=-, leftmargin=0em, itemindent=3.35em, nosep]
\item Анализ и оценка;
\item Планирование вывода из эксплуатации;
\item Миграция данных;
\item Вывод системы из эксплуатации;
\item Завершение проекта.
\end{itemize}

Цель анализа и оценки -- оценка текущего состояния системы, определение причин завершения жизненного цикла, анализ альтернативных решений (модернизация, замена, вывод из эксплуатации). Действия: аудит системы, анализ производительности, затрат на эксплуатацию, соответствия требованиям, оценка рисков, анализ потребностей бизнеса. Результат: отчет об~оценке системы, рекомендации по дальнейшим действиям.

Цель планирования вывода из эксплуатации -- разработка плана вывода системы из эксплуатации, включающего сроки, ресурсы, этапы, ответственных лиц и процедуры. Действия: определение стратегии вывода из эксплуатации (постепенный, одномоментный), разработка графика, определение ресурсов, подготовка документации, уведомление пользователей. Результат: план вывода системы из эксплуатации.

Цель миграции данных: перенос данных из старой системы в новую систему или архив. Действия: анализ данных, очистка данных, разработка процедур миграции, тестирование миграции, фактическая миграция данных. Результат: успешно перенесенные данные, обеспечение целостности и~достоверности данных.

Цель вывода системы из эксплуатации: отключение системы, утилизация оборудования. Действия: остановка работы системы, архивирование данных в соответствии с требованиями, удаление программного обеспечения, утилизация оборудования.

Цель завершения проекта -- формальное завершение проекта по выводу системы из эксплуатации, подведение итогов, анализ полученного опыта. Действия: подготовка итогового отчета, проведение собрания команды проекта, анализ уроков, закрытие контрактов.
\subsection{Оценка эффективности ИС}
Оценка эффективности ИС -- это процесс определения степени достижения целей, поставленных при создании системы. Оценка должна проводиться на протяжении всего жизненного цикла ИС, но особенно важна на~этапе завершения, поскольку позволяет оценить вклад системы в~достижение бизнес-целей и принять обоснованное решение о ее дальнейшей судьбе.

Существуют различные методы оценки эффективности ИС, которые можно разделить на следующие группы:
\begin{enumerate}[leftmargin=0em, itemindent=3.75em, nosep, label=\arabic*)]
\item Экономическая эффективность: оценка затрат, выгод от использования системы. Основные показатели: ROI (Return on Investment), NPV (Net Present Value), IRR (Internal Rate of Return), Payback Period;
\item Функциональная эффективность -- это оценка соответствия функциональности системы требованиям пользователей. Основные показатели: количество поддерживаемых функций, удовлетворенность пользователей;
\item Техническая эффективность: оценка надежности, производительности, масштабируемости и безопасности системы. основные показатели: время безотказной работы, количество ошибок, время отклика системы, количество обрабатываемых транзакций в единицу времени;
\item Организационная эффективность -- оценка влияния системы на~организационную структуру, бизнес-процессы и эффективность работы персонала. Основные показатели: снижение затрат на персонал, ускорение бизнес-процессов, улучшение качества обслуживания клиентов.
\end{enumerate}

Методы оценки:
\begin{enumerate}[leftmargin=0em, itemindent=3.75em, nosep, label=\arabic*)]
\item Анализ затрат и выгод (Cost-Benefit Analysis): Сравнение всех затрат, связанных с разработкой, внедрением и эксплуатацией ИС, с ожидаемыми доходами и выгодами;
\item Экспертная оценка: Привлечение экспертов для оценки различных аспектов работы ИС;
\item Опросы и анкетирование пользователей: Сбор информации о степени удовлетворенности пользователей работой системы, ее функциональности, удобству использования;
\item Метрики производительности: Сбор и анализ метрик, отражающих производительность системы (например, время отклика, количество транзакций в секунду, уровень использования ресурсов);
\item Сравнение с лучшими практиками (Benchmarking): Сравнение характеристик ИС с характеристиками аналогичных систем, реализованных в~других организациях.
\end{enumerate}
\subsection{Утилизация оборудования и архивирование данных}
После вывода системы из эксплуатации необходимо правильно утилизировать оборудование и архивировать данные.

Утилизация оборудования должна осуществляться в соответствии с~законодательством и экологическими нормами. Необходимо обеспечить безопасное удаление данных с жестких дисков и других носителей информации. Оборудование может быть продано, передано другим организациям или~утилизировано.

Архивирование данных необходимо для обеспечения соответствия требованиям законодательства и сохранения возможности доступа к~данным в~будущем. Данные должны быть заархивированы в формате, обеспечивающем их долгосрочное хранение и возможность извлечения. Должны быть разработаны политики хранения и доступа к архивным данным.

Завершение жизненного цикла ИС -- это ответственный процесс, который требует тщательного планирования и подготовки. Правильный вывод системы из эксплуатации позволяет минимизировать риски и издержки, связанные с устаревшей системой, и обеспечить плавный переход к новым технологиям. Оценка эффективности ИС является неотъемлемой частью процесса завершения жизненного цикла, позволяющая получить объективную картину о вкладе системы в достижение бизнес-целей и принять обоснованное решение о ее дальнейшей судьбе.
\newpage
{
\sloppy
\nocite{*}
\cleardoublepage
\addcontentsline{toc}{section}{Список использованных источников}
\bibliographystyle{plainurl}
\bibliography{report2}
}
\end{document}
